%% ============================================================
%%  preamble.tex — Packages, macros, and styling
%%  Comprehensive Code Audit Report template.
%%  Edit the placeholders marked <EDIT> below, then leave the
%%  rest of this file alone.
%% ============================================================

%% ----- Report-wide placeholders ------------------------------------------
% <EDIT> Set these once, at the very top (they're used by \hypersetup below,
% which expands them immediately, so they must be defined before that point).
\newcommand{\reportsystem}{Example Application}
\newcommand{\reportdate}{2030-01-01}

%% ----- Encoding & fonts -------------------------------------------------------
\usepackage[T1]{fontenc}
\usepackage[utf8]{inputenc}
\usepackage{lmodern}
\usepackage{microtype}
\usepackage{amssymb}
\usepackage{textcomp}

%% ----- Page geometry ----------------------------------------------------------
\usepackage[a4paper, top=3.2cm, bottom=3cm, left=2.8cm, right=2.8cm, headheight=14pt]{geometry}

%% ----- Graphics ---------------------------------------------------------------
\usepackage{graphicx}

%% ----- Colour definitions -----------------------------------------------------
\usepackage[dvipsnames,svgnames,table]{xcolor}

\definecolor{sevCritical}{RGB}{180,0,0}
\definecolor{sevHigh}{RGB}{200,80,0}
\definecolor{sevMedium}{RGB}{170,120,0}
\definecolor{sevLow}{RGB}{30,130,30}
\definecolor{sevInfo}{RGB}{50,80,160}
\definecolor{bgCritical}{RGB}{255,240,240}
\definecolor{bgHigh}{RGB}{255,245,235}
\definecolor{bgMedium}{RGB}{255,252,220}
\definecolor{bgLow}{RGB}{240,255,240}
\definecolor{bgInfo}{RGB}{235,240,255}

\definecolor{codebg}{RGB}{248,248,248}
\definecolor{codeframe}{RGB}{200,200,200}
\definecolor{rulegrey}{RGB}{120,120,120}
\definecolor{titlegrey}{RGB}{50,50,70}
\definecolor{linkblue}{RGB}{0,70,150}

%% ----- Hyperlinks -------------------------------------------------------------
% <EDIT> Replace the pdftitle/pdfauthor/pdfsubject below.
\usepackage[
  colorlinks=true,
  linkcolor=linkblue,
  urlcolor=linkblue,
  citecolor=linkblue,
  bookmarks=true,
  bookmarksnumbered=true,
  pdftitle={\reportsystem\ Security Assessment Report},
  pdfauthor={Static Security Review --- \reportdate},
  pdfsubject={White-box Security Assessment},
]{hyperref}
\usepackage{url}
\urlstyle{tt}

%% ----- Tables ----------------------------------------------------------------
\usepackage{booktabs}
\usepackage{longtable}
\usepackage{tabularx}
\usepackage{array}
\usepackage{multirow}
\usepackage{makecell}
\usepackage{colortbl}
\renewcommand{\arraystretch}{1.35}

%% ----- Lists -----------------------------------------------------------------
\usepackage{enumitem}
\setlist[itemize]{noitemsep, topsep=4pt, leftmargin=1.8em}
\setlist[enumerate]{noitemsep, topsep=4pt, leftmargin=2em}

%% ----- Paragraph spacing -----------------------------------------------------
\usepackage[parfill]{parskip}
\setlength{\parindent}{0pt}

%% ----- Code listings ---------------------------------------------------------
\usepackage{listings}

\lstdefinestyle{basestyle}{
  backgroundcolor=\color{codebg},
  frame=single,
  rulecolor=\color{codeframe},
  basicstyle=\footnotesize\ttfamily,
  breaklines=true,
  breakatwhitespace=false,
  showstringspaces=false,
  tabsize=2,
  captionpos=b,
  numbers=none,
  xleftmargin=6pt,
  xrightmargin=6pt,
  framexleftmargin=6pt,
  framexrightmargin=6pt,
  framesep=4pt,
  extendedchars=true,
  inputencoding=utf8,
}

\lstdefinestyle{java}{
  style=basestyle,
  language=Java,
  keywordstyle=\color{NavyBlue}\bfseries,
  stringstyle=\color{ForestGreen},
  commentstyle=\color{gray}\itshape,
}

\lstdefinestyle{xml}{
  style=basestyle,
  language=XML,
  keywordstyle=\color{NavyBlue}\bfseries,
  stringstyle=\color{ForestGreen},
  commentstyle=\color{gray}\itshape,
  morestring=[b]",
  morecomment=[s]{<!--}{-->},
}

\lstdefinestyle{sql}{
  style=basestyle,
  language=SQL,
  keywordstyle=\color{NavyBlue}\bfseries,
  stringstyle=\color{ForestGreen},
  commentstyle=\color{gray}\itshape,
  morekeywords={UPDATE,INSERT,SELECT,FROM,WHERE,LIKE,UNION,HAVING,ORDER,BY,LIMIT,ASC,DESC,SET,AND,OR,NOT},
}

\lstdefinestyle{http}{
  style=basestyle,
  language={},
  commentstyle=\color{gray}\itshape,
  morecomment=[l]{\#},
}

\lstdefinestyle{bash}{
  style=basestyle,
  language=bash,
  keywordstyle=\color{NavyBlue}\bfseries,
  stringstyle=\color{ForestGreen},
  commentstyle=\color{gray}\itshape,
}

\lstdefinestyle{properties}{
  style=basestyle,
  language={},
  morecomment=[l]{\#},
  commentstyle=\color{gray}\itshape,
  stringstyle=\color{ForestGreen},
}

%% ----- TcolorBox for finding blocks ------------------------------------------
\usepackage[most,skins,breakable,listings]{tcolorbox}

%% Severity badge command
\newcommand{\sevbadge}[2]{%
  \colorbox{#1}{\color{white}\bfseries\footnotesize\strut\ #2\ }%
}
\newcommand{\Critical}{\sevbadge{sevCritical}{CRITICAL}}
\newcommand{\High}{\sevbadge{sevHigh}{HIGH}}
\newcommand{\Medium}{\sevbadge{sevMedium}{MEDIUM}}
\newcommand{\Low}{\sevbadge{sevLow}{LOW}}
\newcommand{\Info}{\sevbadge{sevInfo}{INFO}}

\newcommand{\Confirmed}{\textcolor{ForestGreen}{\bfseries Confirmed}}
\newcommand{\Suspected}{\textcolor{Goldenrod}{\bfseries Suspected}}

%% Finding environments by severity
\tcbset{
  findingbase/.style={
    breakable,
    enhanced,
    colback=white,
    boxrule=1pt,
    arc=2pt,
    left=8pt, right=8pt, top=6pt, bottom=6pt,
    attach boxed title to top left={yshift=-2mm, xshift=8pt},
    boxed title style={
      arc=2pt,
      boxrule=0pt,
      left=4pt, right=4pt,
    },
    fonttitle=\bfseries\small,
    separator sign={\ —\ },
  }
}

\newtcolorbox{findingcritical}[2]{
  findingbase,
  colframe=sevCritical,
  colbacktitle=sevCritical,
  coltitle=white,
  title={#1 \textnormal{---} #2},
}
\newtcolorbox{findinghigh}[2]{
  findingbase,
  colframe=sevHigh,
  colbacktitle=sevHigh,
  coltitle=white,
  title={#1 \textnormal{---} #2},
}
\newtcolorbox{findingmedium}[2]{
  findingbase,
  colframe=sevMedium,
  colbacktitle=sevMedium,
  coltitle=white,
  title={#1 \textnormal{---} #2},
}
\newtcolorbox{findinglow}[2]{
  findingbase,
  colframe=sevLow,
  colbacktitle=sevLow,
  coltitle=white,
  title={#1 \textnormal{---} #2},
}
\newtcolorbox{findinginfo}[2]{
  findingbase,
  colframe=sevInfo,
  colbacktitle=sevInfo,
  coltitle=white,
  title={#1 \textnormal{---} #2},
}

%% Convenience macros inside finding boxes
\newcommand{\findsection}[1]{\vspace{4pt}\textbf{\textit{#1:}}\quad}
\newcommand{\findmeta}[2]{\textbf{#1:} #2\enspace}

%% Note / warning box
\newtcolorbox{warningbox}{
  breakable,
  colback=yellow!8!white,
  colframe=Goldenrod!80,
  left=6pt, right=6pt, top=4pt, bottom=4pt,
  arc=2pt,
  boxrule=0.8pt,
}

\newtcolorbox{notebox}{
  breakable,
  colback=blue!5!white,
  colframe=linkblue!60,
  left=6pt, right=6pt, top=4pt, bottom=4pt,
  arc=2pt,
  boxrule=0.8pt,
}

%% ----- Headers & footers -----------------------------------------------------
% <EDIT> Replace the system name below.
\usepackage{fancyhdr}
\pagestyle{fancy}
\fancyhf{}
\fancyhead[L]{\small\textcolor{rulegrey}{\reportsystem\ --- Security Assessment}}
\fancyhead[R]{\small\textcolor{rulegrey}{\leftmark}}
\fancyfoot[C]{\small\textcolor{rulegrey}{\thepage}}
\renewcommand{\headrulewidth}{0.4pt}
\renewcommand{\footrulewidth}{0pt}

%% ----- Chapter / section styling ---------------------------------------------
\usepackage{titlesec}
\titleformat{\chapter}[display]
  {\normalfont\Large\bfseries\color{titlegrey}}
  {\chaptertitlename\ \thechapter}{12pt}{\Huge}
\titlespacing*{\chapter}{0pt}{30pt}{20pt}

\titleformat{\section}
  {\normalfont\large\bfseries\color{titlegrey}}
  {\thesection}{1em}{}
\titlespacing*{\section}{0pt}{18pt}{8pt}

\titleformat{\subsection}
  {\normalfont\normalsize\bfseries\color{titlegrey}}
  {\thesubsection}{1em}{}
\titlespacing*{\subsection}{0pt}{12pt}{4pt}

%% ----- Miscellaneous helpers -------------------------------------------------
\usepackage{rotating}
\usepackage{pdflscape}

%% Inline code helper
\newcommand{\code}[1]{\texttt{#1}}
\newcommand{\file}[1]{\texttt{#1}}
\newcommand{\cls}[1]{\texttt{#1}}
\newcommand{\ep}[1]{\texttt{#1}}   %% endpoint
\newcommand{\fref}[1]{\hyperref[#1]{\textbf{#1}}}  %% finding reference

%% Severity in tables
\newcommand{\tabCrit}{\cellcolor{sevCritical!20}\textcolor{sevCritical}{\bfseries Critical}}
\newcommand{\tabHigh}{\cellcolor{sevHigh!20}\textcolor{sevHigh}{\bfseries High}}
\newcommand{\tabMed}{\cellcolor{sevMedium!20}\textcolor{sevMedium}{\bfseries Medium}}
\newcommand{\tabLow}{\cellcolor{sevLow!20}\textcolor{sevLow}{\bfseries Low}}
\newcommand{\tabInfo}{\cellcolor{sevInfo!20}\textcolor{sevInfo}{\bfseries Info}}

%% Placeholder marker -- wrap text that needs to be filled in manually
\newcommand{\fillme}[1]{\colorbox{yellow!40}{\ttfamily\bfseries\small [ FILL: #1 ]}}
